\usepackage[T1]{fontenc}
\usepackage[utf8]{inputenc}
\usepackage{lmodern}
\usepackage[margin=1in]{geometry}
\usepackage{amsmath,amssymb,amsthm,mathtools}
% Inline mathematics keeps normal text style; displays use display style.
\usepackage{enumitem}
\usepackage{xcolor}
\usepackage{tikz,tikz-cd}
\usetikzlibrary{arrows.meta,positioning,calc}
\usepackage{hyperref}
\usepackage[nameinlink,capitalize,noabbrev]{cleveref}
\hypersetup{
  colorlinks=true,
  linkcolor=blue!50!black,
  citecolor=blue!50!black,
  urlcolor=blue!50!black,
  hypertexnames=false,
  pdftitle={Lecture Notes on Mathematical Analysis: Single and Multivariable Calculus},
  pdfauthor={Tommy Chen},
  pdfsubject={Version 3.30.0: mathematical analysis, elementary functions, and structural appendices},
  pdfkeywords={mathematical analysis, calculus, real analysis, multivariable calculus, Stokes theorem}
}
\newcommand{\notesversion}{Version 3.30.0}
\newcommand{\notesdate}{September 19, 2026}
\newcommand{\noteslicense}{%
  This work is licensed under the
  \href{https://creativecommons.org/licenses/by/4.0/}
  {Creative Commons Attribution 4.0 International License (CC BY 4.0).}%
}
\setlist[enumerate]{label=(\alph*),itemsep=0.35em}
% Number chapters and sections only; retain unnumbered subsections in the contents.
\setcounter{secnumdepth}{1}
% Give paragraph breaking room for the longer section-based reference labels.
\setlength{\emergencystretch}{2em}
\usepackage{needspace}
\usepackage{placeins}
\usepackage{float}
\setcounter{tocdepth}{2}
% Allow two-digit chapter and section numbers such as 13.10 in the contents.
\makeatletter
\renewcommand*\l@section{\@dottedtocline{1}{1.5em}{3em}}
\renewcommand*\l@subsection{\@dottedtocline{2}{4.5em}{3.8em}}
\makeatother
\newtheoremstyle{analysisplain}
  {\topsep}
  {\topsep}
  {\upshape}
  {}
  {\bfseries}
  {.}
  {.5em}
  {}
\theoremstyle{analysisplain}
\newtheorem{theorem}{Theorem}[section]
\newtheorem{proposition}[theorem]{Proposition}
\newtheorem{lemma}[theorem]{Lemma}
\newtheorem{corollary}[theorem]{Corollary}
\newtheorem{remark}[theorem]{Remark}
\theoremstyle{definition}
\newtheorem{definition}[theorem]{Definition}
\newtheorem{construction}[theorem]{Construction}
\newtheorem{example}[theorem]{Example}
\newtheorem{exercise}{Exercise}[chapter]
% Exercises have their own sequence, restarting in each chapter or appendix.
\renewcommand{\theexercise}{\arabic{exercise}}
% Reserve the heading and opening lines before hyperref creates the anchor.
% Otherwise a theorem pushed to the next page can leave its destination behind.
\BeforeBeginEnvironment{theorem}{\par\Needspace{6\baselineskip}}
\BeforeBeginEnvironment{proposition}{\par\Needspace{6\baselineskip}}
\BeforeBeginEnvironment{lemma}{\par\Needspace{6\baselineskip}}
\BeforeBeginEnvironment{corollary}{\par\Needspace{6\baselineskip}}
\BeforeBeginEnvironment{definition}{\par\Needspace{6\baselineskip}}
\BeforeBeginEnvironment{construction}{\par\Needspace{6\baselineskip}}
\BeforeBeginEnvironment{example}{\par\Needspace{6\baselineskip}}
\BeforeBeginEnvironment{remark}{\par\Needspace{6\baselineskip}}
\BeforeBeginEnvironment{exercise}{\par\Needspace{6\baselineskip}}
\AtBeginEnvironment{definition}{\let\emph\textbf}
\crefname{construction}{construction}{constructions}
\Crefname{construction}{Construction}{Constructions}
\newcommand{\R}{\mathbb{R}}
\newcommand{\Q}{\mathbb{Q}}
\newcommand{\N}{\mathbb{N}}
\newcommand{\Z}{\mathbb{Z}}
\newcommand{\C}{\mathbb{C}}
\newcommand{\eps}{\varepsilon}
\newcommand{\abs}[1]{\lvert #1\rvert}
\newcommand{\norm}[1]{\lVert #1\rVert}
\newcommand{\set}[1]{\left\{#1\right\}}
\newcommand{\dd}{\mathop{}\!\mathrm{d}}

% Use complete appendix titles in both the TOC and the PDF outline.
% Chapter counters and destinations still come from the ordinary chapter command.
\newcommand{\appendixchapter}[1]{%
  \let\appendixsavedcontentsline\addcontentsline
  \renewcommand{\addcontentsline}[3]{%
    \appendixsavedcontentsline{##1}{##2}{Appendix \thechapter: #1}}%
  \chapter{#1}%
  \let\addcontentsline\appendixsavedcontentsline
}

% Hinted exercises share the ordinary exercise counter and use label destinations.
\newtheoremstyle{hintedanalysis}{\topsep}{\topsep}{\upshape}{}{\bfseries}{}{.5em}
  {\thmname{#1}\thmnumber{ #2}.\textsuperscript{\hyperref[hint:\currenthint]{*}}}
\theoremstyle{hintedanalysis}
\newtheorem{hintedexercise}[exercise]{Exercise}
\newenvironment{hexercise}[1]
  {\def\currenthint{#1}\begin{hintedexercise}\label{ex:#1}}
  {\end{hintedexercise}}
\BeforeBeginEnvironment{hexercise}{\par\Needspace{6\baselineskip}}
\newcommand{\exercisehint}[2]{%
  \par\Needspace{5\baselineskip}\phantomsection\label{hint:#1}%
  \noindent\textbf{\hyperref[ex:#1]{Exercise \ref*{ex:#1}.}}\quad #2\par\medskip}
